\section{БЛОК 1 (50 баллов)}

\small
\setlength{\parskip}{1pt}
\setlength{\abovedisplayskip}{3pt}
\setlength{\belowdisplayskip}{3pt}
\setlength{\abovedisplayshortskip}{2pt}
\setlength{\belowdisplayshortskip}{2pt}

\begin{multicols}{2}

\begin{definition}
    Дифференциал функции $f(x_1, \ldots, x_n)$:
    \[
        df = f'_{x_1}dx_1 + \ldots + f'_{x_n}dx_n
    \]
\end{definition}

\begin{definition}
    Градиент функции $f(x_1, \ldots, x_n)$:
    \[
        \nabla f =
        \left(
            \frac{\partial f}{\partial x_1},
            \ldots,
            \frac{\partial f}{\partial x_n}
        \right)
    \]
\end{definition}

\begin{definition}
    Производная функции $f$ в точке $M_0$ по направлению вектора $\vec l$:
    \[
        \frac{\partial f}{\partial l}(M_0)
        =
        \left(
            \nabla f\big|_{M_0},
            \frac{\vec l}{|\vec l|}
        \right)
    \]
\end{definition}

\begin{definition}
    Критические точки функции $f(x_1, \ldots, x_n)$ находятся из системы:
    \[
        \begin{cases}
            f'_{x_1} = 0, \\
            f'_{x_2} = 0, \\
            \vdots \\
            f'_{x_n} = 0.
        \end{cases}
    \]
\end{definition}

\begin{theorem}[Необходимое условие экстремума]
    \[
        x^0 \text{ --- точка экстремума функции } f
        \Rightarrow
    \]
    \[
        x^0 \text{ --- критическая точка.}
    \]
\end{theorem}

\begin{definition}
    Матрица Гессе функции $f(x_1, \ldots, x_n)$:
    \[
        H =
        \begin{pmatrix}
            f''_{x_1x_1} & \ldots & f''_{x_1x_n} \\
            \vdots & \ddots & \vdots \\
            f''_{x_nx_1} & \ldots & f''_{x_nx_n}
        \end{pmatrix}
    \]
\end{definition}

\begin{remark}
    Алгоритм исследования на экстремум:
    \[
    \begin{array}{ll}
        1) & \text{найти первые частные производные;} \\
        2) & \text{найти критические точки;} \\
        3) & \text{найти вторые частные производные;} \\
        4) & \text{составить матрицу Гессе;} \\
        5) & \text{проверить по критерию Сильвестра.}
    \end{array}
    \]
\end{remark}

\begin{theorem}[Критерий Сильвестра для экстремума]
    Пусть $H$ --- матрица Гессе в критической точке.

    Если форма положительно определена:
    \[
        \Delta_1 > 0,\ \Delta_2 > 0,\ \ldots,\ \Delta_n > 0
        \Rightarrow \min
    \]

    Если форма отрицательно определена:
    \[
        \Delta_1 < 0,\ \Delta_2 > 0,\ \Delta_3 < 0,\ \ldots
        \Rightarrow \max
    \]

    Если форма не положительно и не отрицательно определена, то точка не является точкой экстремума.
\end{theorem}

\begin{definition}
    Если $u = u(x,y)$ задана неявно:
    \[
        F(x,y,u) = 0,
    \]
    то
    \[
        u'_x = -\frac{F'_x}{F'_u},
        \qquad
        u'_y = -\frac{F'_y}{F'_u}.
    \]
\end{definition}

\begin{remark}
    В критических точках $u'_x = u'_y = 0$, поэтому для неявно заданных функций удобно:
    \[
        u''_{xx} = -\frac{F''_{xx}}{F'_u},\quad
        u''_{yy} = -\frac{F''_{yy}}{F'_u},\quad
        u''_{xy} = -\frac{F''_{xy}}{F'_u}.
    \]
\end{remark}

\end{multicols}




\section{БЛОК 2 (70 баллов)}

\small
\setlength{\parskip}{1pt}
\setlength{\abovedisplayskip}{3pt}
\setlength{\belowdisplayskip}{3pt}
\setlength{\abovedisplayshortskip}{2pt}
\setlength{\belowdisplayshortskip}{2pt}

\begin{multicols}{2}

\begin{definition}
    Последовательность функций $\{f_n\}$ сходится к функции $f$ поточечно на $E$, если
    \[
        \forall x \in E
        \quad
        \lim_{n \to \infty} f_n(x) = f(x).
    \]
\end{definition}

\begin{definition}
    Последовательность функций $\{f_n\}$ сходится равномерно к функции $f$ на $E$, если
    \[
        \forall \varepsilon > 0\ \exists N\ \forall n > N\ \forall x \in E:
        |f_n(x) - f(x)| < \varepsilon
    \]
    \[
        \Longleftrightarrow
    \]
    \[
        \lim_{n \to \infty}
        \sup_{x \in E}|f_n(x) - f(x)| = 0.
    \]
\end{definition}

\begin{theorem}
    Равномерная сходимость влечёт поточечную сходимость:
    \[
        f_n \rightrightarrows f \text{ на } E
        \Rightarrow
        f_n \to f \text{ поточечно на } E.
    \]
\end{theorem}

\begin{definition}
    Функциональным рядом называется ряд вида
    \[
        \sum_{n=1}^{\infty} f_n(x),
    \]
    где $f_n(x)$ --- функции, заданные на множестве $E$.
\end{definition}

\begin{definition}
    Степенным рядом называется функциональный ряд вида
    \[
        \sum_{n=0}^{\infty} C_n(x-a)^n.
    \]
\end{definition}

\begin{theorem}[Признак Вейерштрасса]
    \[
        \forall x \in E,\ \forall n:\ |f_n(x)| \leq a_n
    \]
    и
    $
        \sum_{n=1}^{\infty} a_n \text{ сходится},
    $
    $\Rightarrow$
    $\sum_{n=1}^{\infty} f_n(x)$
    сходится равномерно на $E$.
\end{theorem}

\begin{theorem}[Радиус сходимости]
    Ряд
    \[
        \sum_{n=0}^{\infty} C_n(x-a)^n
    \]
    сходится при $|x-a|<R$, расходится при $|x-a|>R$.
    Концы $a-R$ и $a+R$ проверяются отдельно. То есть подставляем эти концы и решаем задачу как с КР4.
    \[
        \frac{1}{R}=\lim_{n\to\infty}\sqrt[n]{|C_n|},
        \qquad
        R=\lim_{n\to\infty}\left|\frac{C_n}{C_{n+1}}\right|.
    \]
\end{theorem}

\begin{lemma}[Область сходимости степенного ряда]
        Область сходимости --- это множество всех $x$, при которых ряд сходится.

    \[
        |x-a|<R
        \Longleftrightarrow
        a-R<x<a+R.
    \]
    Концы $a-R$ и $a+R$ проверяются отдельно.
\end{lemma}

\begin{definition}
    Частная производная функции $f(x,y)$ по $x$ в точке $(x_0, y_0)$:
    \[
        f'_x(x_0,y_0)
        =
        \lim_{\Delta x \to 0}
        \frac{
            f(x_0 + \Delta x, y_0) - f(x_0, y_0)
        }{\Delta x}.
    \]
    
    Аналогично по $y$:
    \[
        f'_y(x_0,y_0)
        =
        \lim_{\Delta y \to 0}
        \frac{
            f(x_0, y_0 + \Delta y) - f(x_0, y_0)
        }{\Delta y}.
    \]
\end{definition}

\begin{definition}
    $f$ дифференцируема в $(x_0,y_0)$, если $\exists A,B$:
    \[
        \begin{aligned}
            \Delta f
            &=
            A\Delta x+B\Delta y \\
            &\quad+
            o\left(\sqrt{\Delta x^2+\Delta y^2}\right),
        \end{aligned}
    \]
    где $\Delta f=f(x,y)-f(x_0,y_0)$, $\Delta x=x-x_0$, $\Delta y=y-y_0$.
\end{definition}

\begin{remark}[Необходимое условие дифф. в точке]
    \[
        f \text{ дифф. в } (x_0,y_0)
        \Rightarrow
        \exists f'_x(x_0,y_0),\ \exists f'_y(x_0,y_0)
    \]
    \[
        A = f'_x(x_0,y_0),
        \qquad
        B = f'_y(x_0,y_0).
    \]
\end{remark}

\begin{theorem}[Достаточное условие дифференцируемости]
    Все частные производные функции $f$ непрерывны в некоторой окрестности точки $(x_0,y_0)$, $\Rightarrow$ функция $f$ дифференцируема в этой точке.
\end{theorem}

\begin{definition}
    Ряд Тейлора функции $f$ в точке $a$:
    \[
        \sum_{n=0}^{\infty}
        \frac{f^{(n)}(a)}{n!}(x-a)^n.
    \]
\end{definition}


\begin{definition}
    Ряд Маклорена --- это ряд Тейлора при $a = 0$:
    \[
        f(x)
        =
        \sum_{n=0}^{\infty}
        \frac{f^{(n)}(0)}{n!}x^n.
    \]
\end{definition}

\begin{remark}
    Ряд Маклорена является степенным рядом, поэтому для него ищется радиус сходимости как обычно (выше написано).
    (Для понимания) Иерархия рядов:
    \[
        \text{Ряд Маклорена}
        \subset
        \text{Ряд Тейлора} \subset
    \]
    \[
        \text{Степенной ряд}
        \subset
        \text{Функциональный ряд}.
    \]
\end{remark}

\begin{hint}[Основные ряды Маклорена:]
    \[
        e^x = \sum_{n=0}^{\infty} \frac{x^n}{n!}
    \]
    \[
        \sin x = \sum_{n=0}^{\infty} (-1)^n \frac{x^{2n+1}}{(2n+1)!}
    \]
    \[
        \cos x = \sum_{n=0}^{\infty} (-1)^n \frac{x^{2n}}{(2n)!}
    \]
    \[
        (1+x)^{\alpha}
        =
        \sum_{n=0}^{\infty}
        C_{\alpha}^{n}x^n
    \]
    \[
        \ln(1+x)
        =
        \sum_{n=1}^{\infty}
        (-1)^{n-1}\frac{x^n}{n}
    \]
    \[
        \ln(1-x)
        =
        -\sum_{n=1}^{\infty}
        \frac{x^n}{n}
    \]
    \[
        \arctg x
        =
        \sum_{n=0}^{\infty}
        (-1)^n\frac{x^{2n+1}}{2n+1}
    \]
    \[
        \arcsin x
        =
        \sum_{n=0}^{\infty}
        \frac{(2n)!}{4^n(n!)^2(2n+1)}x^{2n+1}
    \]
    \[
        \arccos x = \frac{\pi}{2} - \arcsin x
    \]
    \[
        \frac{1}{1-x}
        =
        \sum_{n=0}^{\infty}x^n
    \]
    \[
        \frac{1}{1+x}
        =
        \sum_{n=0}^{\infty}(-1)^n x^n
    \]
    \[
        (1+x)^{\alpha}
        =
        1 + \sum_{n=1}^{\infty} C_{\alpha}^{n}x^n
    \]
    \[
        C_{\alpha}^{n}
        =
        \frac{\alpha(\alpha-1)\ldots(\alpha-n+1)}{n!}.
    \]


\end{hint}


\end{multicols}